All major published papers on reinforcement learning with execution feedback (RLEF) for code LLMs use binary pass/fail reward --- yet under Group Relative Policy Optimization (GRPO), this choice produces zero gradient contribution in 98.7\% of early-training groups when the model cannot produce complete solutions. We analyze the group-level advantage computation in GRPO and prove that ratio reward ($k/n$ test cases passing) mathematically guarantees non-zero advantage variance whenever any completion in a group achieves partial credit --- precisely the regime where binary reward guarantees zero variance. For a realistic 8-completion group with test-pass counts $[0, 0, 1, 2, 0, 3, 0, 0]$, binary reward yields advantage variance 0.0000 while ratio reward yields 0.0475; a 1,000-group simulation confirms this holds for 98.7\% of early-training groups under realistic partial-pass distributions. We further show that observing policy-level effects of this mechanistic difference requires the base model to achieve non-zero partial solve rates on training data: a controlled 208-step GRPO experiment (DeepSeek-Coder-6.7B on APPS, \texttt{max\_new\_tokens=512}) produced \texttt{reward\_mean = 0} and \texttt{clipped\_ratio = 1.0} throughout for both conditions, confirming that 512-token generation is insufficient for syntactically complete APPS solutions and thus collapses both rewards to zero. This precision null result identifies the experimental prerequisite --- non-zero training-set solve rate --- required before ratio reward's mechanistic advantage can manifest as policy-level performance differences. We release a complete, validated GRPO + ratio reward training framework (35/35 unit tests passing) to enable follow-up experiments under correctly-powered conditions.
